La primera respuesta de la IA generó un diagrama \texttt{sequenceDiagram} sintácticamente compilable, pero con múltiples divergencias arquitectónicas respecto al glosario común y al modelo estructural de datos:
\begin{itemize}[leftmargin=*,nosep]
\item \textbf{Identificadores de participantes externos incorrectos:} La IA declaró \texttt{participant GW as Pasarela Bizum} y \texttt{participant BK as App Banco}. Se corrigieron para respetar estrictamente el glosario común: \texttt{participant PB as Pasarela Bizum} y \texttt{participant BK as App del banco}.
\item \textbf{Invención de cuatro endpoints y estados ajenos al modelo ER:} La IA inventó \texttt{POST /api/orders/checkout/bizum}, \texttt{/api/v1/bizum/init}, \texttt{/api/webhooks/bizum} y \texttt{GET /api/orders/\{id\}/status}, así como los estados \texttt{PAGADO}, \texttt{PAGO\_FALLIDO}, \texttt{COMPLETADO} y \texttt{CANCELADO}. Se unificó la entrada con el endpoint oficial \texttt{POST /api/pedidos} y los estados válidos de \texttt{PEDIDO} (\texttt{PENDIENTE}) y \texttt{PAGO} (\texttt{APROBADO} o \texttt{RECHAZADO}).
\item \textbf{Trazabilidad con las clases \texttt{PasarelaPago}, \texttt{PagoBizum} y \texttt{StockTalla}:} Se formalizaron los mensajes técnicos mediante llamadas a métodos: \texttt{validarFormatoTelefono(telefono)}, \texttt{StockTalla.reservar(cantidad)}, \texttt{savePedido(PENDIENTE)}, \texttt{PasarelaPago.cobrar(PagoBizum)}, \texttt{PagoBizum.solicitarAutorizacion()}, \texttt{savePago(APROBADO)}, \texttt{descontarStock()}, \texttt{vaciarCarrito()} y \texttt{liberarReservaStock()}, junto con la directiva \texttt{mirrorActors: false}.
\end{itemize}
